% Source r6 block, retained for companion integration.
\section{Bounded rise cannot remain coprime to fresh moduli}\label{sec:barrier}

The arithmetic obstruction needed for Theorem~\ref{res:bounded} is a
first-crossing argument: a numerator with bounded upward increments cannot
avoid a persistent block of forbidden heights.

\begin{lemma}[shifted blocks of consecutive multiples]\label{res:crt}
Let $m_0,\ldots,m_{B-1}$ be pairwise coprime and at least $2$.  For every bound
there is a $t$ beyond it with $m_i\mid t+i$ for each $i<B$.
\end{lemma}

\begin{proof}
Standard Chinese remaindering: solve $t\equiv-i \pmod{m_i}$ simultaneously,
then add multiples of $\prod_i m_i$ to pass the bound.
\end{proof}

\noindent Formalised as the
\lword{Erdos243/ReciprocalTailRigidity.lean}{833}{exists_shifted_consecutiveMultiples}{shifted
consecutive multiples}.

\begin{example}[a block of three consecutive multiples]\label{ex:crt}
Take $B=3$ and $(m_0,m_1,m_2)=(3,4,5)$.  The congruences $t\equiv0\pmod3$,
$t\equiv3\pmod4$ and $t\equiv3\pmod5$ hold exactly when $t\equiv3\pmod{60}$.
At $t=3$ the block is $3,4,5$ with $3\mid3$, $4\mid4$ and $5\mid5$; at $t=63$
it is $63,64,65$ with $3\mid63$, $4\mid64$ and $5\mid65$.  A sequence of
natural numbers that starts at $0$, tends to infinity and rises by at most $3$
at each step cannot step over the block $\{3,4,5\}$: it has a first value at
least $3$, that value is at most $5$, and every member of $\{3,4,5\}$ is
divisible by one of the three moduli.
\end{example}

\begin{theorem}[bounded rise cannot remain coprime to fresh moduli]\label{res:barrier}
Let $u:\N\to\N$ tend to infinity with $u_{n+1}\le u_n+B$ for a fixed $B>0$.
Then $u$ cannot remain coprime to infinitely many fresh pairwise coprime
moduli: there is no family of pairwise coprime $m_i\ge2$, one for each index,
such that $\gcd(m_i,u_t)=1$ whenever $i<t$.
\end{theorem}

\begin{proof}
Take $B$ of the moduli, from indices beyond any chosen point, and use
Lemma~\ref{res:crt} to find $t$ far out with the $i$th of them dividing $t+i$.
The interval $[t,t+B)$ has length $B$, and $u$ climbs by at most $B$ per step
while tending to infinity, so some $u_s$ lands in it.  Then $u_s=t+i$ for some
$i<B$, and the $i$th modulus divides $u_s$.  Since that modulus is at least
$2$, this contradicts $\gcd(m_i,u_s)=1$.
\end{proof}

\noindent Formalised as the
\lword{Erdos243/ReciprocalTailRigidity.lean}{897}{no_boundedRise_of_tailAvoidance}{bounded-rise
barrier}.  The key point is that the two hypotheses pull against each other
exactly: divergence forces $u$ to travel arbitrarily far, and bounded rise
forbids it from stepping over a window of width $B$.  Uniformity of the bound
is what the proof uses, and it is the hypothesis that cannot be relaxed by this
argument: a rise bounded only along a subsequence leaves the intervening steps
free to jump the window.  The variable-rise construction discussed in
Problem~\ref{res:variablerise} now supplies such sequences, even under
arbitrarily slowly growing unbounded envelopes.  The fixed-bound argument
therefore needs a joint growth-and-activation hypothesis before it can be
extended.

The moduli are indexed by the same set as the sequence, and the $i$th of them
is asked to divide no value $u_t$ with $t>i$; nothing at all is asked of it at
or before its own index.  That asymmetry is what makes the hypothesis available
in the application below, where $m_i$ is the $i$th multiplier and only the
later numerators are known to be coprime to it.

The barrier applies to the problem because a \emph{reduced} exact tail already
carries pairwise coprime moduli.  Call $(u,v)$ a reduced exact tail with
multipliers $a$ when $\gcd(u_n,v_n)=1$,
\[
 u_{n+1}+v_n=a_nu_n,\qquad v_{n+1}=a_nv_n .
\]
Here $u$ is the numerator and $v$ the denominator: the two recurrences are
those of Section~\ref{sec:state}, with $u$ in the role of the tail state $C$
and $v$ in that of the denominator state $D$, and with coprimality imposed at
every index.

\begin{proposition}[reduced tails are pairwise coprime]\label{res:reduced}
In a reduced exact tail, $a_n$ is coprime to $v_n$; the multipliers at distinct
indices are pairwise coprime; and every earlier multiplier is coprime to every
later numerator.
\end{proposition}

\noindent These are the
\lword{Erdos243/ReciprocalTailRigidity.lean}{983}{reducedStep_coprime_currentFactor}{step
coprimality}, the
\lword{Erdos243/ReciprocalTailRigidity.lean}{1010}{reducedTail_pairwiseCoprime}{pairwise
coprimality}, and the
\lword{Erdos243/ReciprocalTailRigidity.lean}{1024}{reducedTail_wholeModulusAvoidance}{whole-modulus
avoidance}.  Feeding them to Theorem~\ref{res:barrier} gives the form used
later: there is no reduced exact tail whose numerator tends to infinity with a
uniformly bounded upward increment, the
\lword{Erdos243/ReciprocalTailRigidity.lean}{1035}{no_boundedRise_reducedTail}{reduced
bounded-rise exclusion}, together with its eventual form, the
\lword{Erdos243/ReciprocalTailRigidity.lean}{1060}{no_eventuallyBoundedRise_reducedTail}{eventual
reduced exclusion}.

\begin{example}[Sylvester's sequence as a reduced exact tail]\label{ex:reduced}
Put $u_n=1$ and $v_n=D_n=1,2,6,42,1806,\ldots$ as in
Example~\ref{ex:sylvester}, with multipliers $a_n=v_n+1=2,3,7,43,1807,\ldots$.
Then $\gcd(u_n,v_n)=1$, $u_{n+1}+v_n=1+v_n=a_nu_n$ and $v_{n+1}=a_nv_n$, so
this is a reduced exact tail, and Proposition~\ref{res:reduced} returns the
classical fact that Sylvester's numbers are pairwise coprime.  Its numerator is
constant, so it satisfies every hypothesis of the exclusion just stated except
divergence; since the tail exists, that hypothesis cannot be dropped.
\end{example}

An arbitrary orbit of Section~\ref{sec:state} need not be reduced, so one more step is needed
before the barrier can be applied to it: the common factor must stop changing.
We call a set of indices \emph{cofinal} when it is infinite, and say that a
property holds \emph{cofinally} when the set of indices at which it holds is
cofinal.  The step below is the only one connecting the orbit of
Section~\ref{sec:state} to the reduced tails of
Proposition~\ref{res:reduced}, and it is where the cofinal boundedness in the
proof of Theorem~\ref{res:bounded} is used.

\begin{proposition}[the tail gcd stabilises]\label{res:gcdstab}
Let $(a,D,C)$ be an exact natural orbit, that is, a triple of $\N$-valued
sequences with $C_{n+1}+D_n=a_nC_n$ and $D_{n+1}=a_nD_n$, whose centred state
$E_n=\ctr(a_n,D_n,C_n)$ is negative cofinally, with magnitudes bounded along
that cofinal set.  Then
$\gcd(C_n,D_n)$ is eventually constant, and beyond that index the orbit divided
by the stable gcd is a reduced exact tail.
\end{proposition}

\begin{proof}
The tail gcd $\gcd(C_n,D_n)$ divides its successor, the
\lword{Erdos243/ReciprocalTailRigidity.lean}{1424}{tailGcd_dvd_succ}{gcd
monotonicity}, so the tail gcds form a positive divisibility chain.  At an
index where the centred state is negative the tail gcd is exactly
$\gcd(C_n,e_n)$, the
\lword{Erdos243/ReciprocalTailRigidity.lean}{1581}{tailGcd_eq_gcd_negativeMagnitude}{gcd
at a negative index}, and is therefore at most the magnitude there; so the
chain is bounded along the cofinal set of negative indices.  A positive
divisibility chain whose values are bounded along a cofinal set is eventually
constant, the
\lword{Erdos243/ReciprocalTailRigidity.lean}{1599}{dvdChain_eventuallyConstant_of_cofinally_bounded}{chain
stabilisation}.  Hence cofinally bounded negative magnitudes make the tail gcd
stabilise, the
\lword{Erdos243/ReciprocalTailRigidity.lean}{1630}{tailGcd_eventuallyConstant_of_cofinally_boundedNegative}{gcd
stabilisation}, and beyond that point dividing by the stable gcd leaves the
numerator and denominator coprime, that is, a reduced exact tail.
\end{proof}

\noindent The stabilisation of the tail gcd is the checked statement cited at
the end of the proof; the last clause, that the divided orbit is a reduced
exact tail, is exposition and carries no separate label.  Boundedness is
required only \emph{along a cofinal set}.  That is what the situation gives:
the tail gcd is
controlled only at the negative indices, and cofinally many of them is all the
divisibility chain needs.

Normalised vanishing by itself gives a different, strictly weaker conclusion.
For an exact natural orbit with $C_n>0$, put
\[
 G_n=\gcd(C_n,D_n),\qquad
 \Gamma(N)=\#\{0\le j<N:G_j<G_{j+1}\}.
\]

\begin{proposition}[normalised vanishing makes strict gcd changes sparse]
\label{res:gcdsparse}
Assume
\[
  (\forall K\ge1)(\exists N)(\forall n\ge N)\qquad
  K|E_n|<C_n.
\]
Then, for every $K\ge1$, eventually $K\Gamma(N)<N$.  Moreover, for every
starting bound $B$ and block length $L$, some $n\ge B$ satisfies
\[
  G_n=G_{n+1}=\cdots=G_{n+L}.
\]
\end{proposition}

\noindent The first assertion is the
\lword{Erdos243/ReciprocalTailRigidity.lean}{2180}{tailGcd_strictGrowthCount_sublinear_of_normalizedVanishes}{sublinear
strict-growth theorem}; the second is the
\lword{Erdos243/ReciprocalTailRigidity.lean}{2196}{tailGcd_exists_arbitrarilyLate_constBlock_of_normalizedVanishes}{arbitrarily
late constant-block theorem}.  The proof first converts normalised vanishing
into subexponential growth of $C_n$, then charges every strict divisibility
increase of $G_n$ against a finite power-of-two budget.  Sparse strict changes
force long finite gaps between them.

The quantifiers matter.  Proposition~\ref{res:gcdstab} uses cofinally bounded
negative magnitudes and yields eventual constancy.  Proposition~\ref{res:gcdsparse}
uses only normalised vanishing and yields arbitrarily late constant blocks of
each prescribed finite length.  It does not yield one infinite constant tail,
and it does not exclude cofinally unbounded negative excursions.
